\documentclass{resume} % Use the custom resume.cls style

\usepackage[left=0.45in,top=0.4in,right=0.45in,bottom=0.4in]{geometry}
\usepackage{enumitem}

% Compact one-to-two page industry CV.
\def\sectionskip{\smallskip}
\def\sectionlineskip{\smallskip}

\newcommand{\tab}[1]{\hspace{.2667\textwidth}\rlap{#1}}
\newcommand{\itab}[1]{\hspace{0em}\rlap{#1}}

\name{Anh Nguyen Hoang}

\address{+84 398 272 620 \\ Hanoi, Vietnam (open to remote)}
\address{
    \href{mailto:anhnh.4work@gmail.com}{anhnh.4work@gmail.com} \\
    \href{https://anhnh2002.github.io/}{Homepage} \\
    \href{https://github.com/anhnh2002}{GitHub}
}

\begin{document}

%----------------------------------------------------------------------------------------
% SUMMARY
%----------------------------------------------------------------------------------------

\begin{rSection}{Summary}

LLM application engineer with two years of hands-on experience shipping
LLM-powered products and autonomous agents, from a consumer AI companion serving
300,000+ users and roughly 100,000 messages per day to an open-source
repository-documentation agent with 1,500+ GitHub stars. Strong background in
agent design (planning, memory, tool calling, multi-agent orchestration),
evaluation frameworks, and agent reliability, backed by publications at ACL,
NAACL, and EMNLP. Comfortable owning a problem end to end, from ambiguous
business need to measured, production-ready AI system.

\end{rSection}

%----------------------------------------------------------------------------------------
% SKILLS
%----------------------------------------------------------------------------------------

\begin{rSection}{Skills}

\begin{tabular}{@{} >{\bfseries}p{0.2\textwidth} @{\hspace{2ex}} p{0.76\textwidth}}
LLM \& Agents &
    Agent architectures (planning, long-term memory, tool calling, multi-agent
    orchestration), prompt engineering and optimization, RAG and vector search,
    Model Context Protocol (MCP) tooling and security, LLM fine-tuning
    (LoRA/full), evaluation harness and benchmark design, LLM-as-judge \\
Engineering &
    Python, PyTorch, Hugging Face Transformers, LLM APIs (OpenAI, Anthropic,
    open-weight models), REST APIs, SQL/NoSQL databases, vector stores, Docker,
    Git, Linux, CI and observability for AI systems \\
Domains &
    Conversational AI, AI for software engineering (code agents, program
    repair, code migration, vulnerability detection), continual learning,
    text-to-speech adaptation \\
\end{tabular}

\end{rSection}

%----------------------------------------------------------------------------------------
% EXPERIENCE
%----------------------------------------------------------------------------------------

\begin{rSection}{Experience}

\textbf{FPT Software AI Center} \hfill May 2025 -- Present\\
\textit{AI Research Resident (LLM agents for software engineering)}
\hfill \textit{Hanoi, Vietnam}

\begin{itemize}[leftmargin=1.5em,itemsep=0em,topsep=0.1em]
    \item Designed and built \textbf{CodeWiki}, an LLM agent framework that
    generates holistic, repository-level documentation for large codebases
    across seven programming languages; open-sourced to 1,500+ GitHub stars
    and published at Findings of ACL 2026
    \item Built the agent's planning, hierarchical decomposition, and
    tool-calling pipeline for repository-scale context, plus an evaluation
    benchmark and LLM-as-judge framework to measure documentation quality
    \item Co-developed \textbf{RustPrint}, a documentation-guided agentic
    system that migrates C codebases to Rust, coordinating multiple
    specialized agents with automated build-and-test feedback loops
    \item Leading \textbf{CHAIR}, a controller that automatically improves agent
    harnesses through hypothesis generation, information-gain experiment
    selection, and controlled single-factor evaluation
    \item Collaborating with UIUC and Mila on agentic vulnerability detection
    and on hardening coding agents against MCP tool-poisoning attacks
\end{itemize}

\vspace{0.3em}

\textbf{Oraichain Labs} \hfill May 2024 -- April 2025\\
\textit{AI Engineer} \hfill \textit{Hanoi, Vietnam}

\begin{itemize}[leftmargin=1.5em,itemsep=0em,topsep=0.1em]
    \item Led development of \textbf{Cupiee}, an emotional AI companion platform
    that grew to 300,000+ users and approximately 100,000 messages per day
    \item Designed the conversational agent stack: persona and prompt design,
    long-term user memory, retrieval over conversation history, and
    safety and behavior guardrails; iterated on prompts and model choices
    against user-engagement and quality metrics
    \item Built autonomous multi-agent systems for social media engagement and
    crypto market analysis, integrating LLMs with external APIs, scheduled
    workflows, and data pipelines
    \item Owned production concerns including latency, cost per message,
    monitoring, and rapid experiment-deploy cycles
\end{itemize}

\vspace{0.3em}

\textbf{Research Assistant, NLP and Continual Learning}
\hfill March 2024 -- April 2025\\
\textit{Advised by Prof.\ Linh Ngo Van (HUST) and Prof.\ Thien Huu Nguyen
(University of Oregon)} \hfill \textit{Hanoi, Vietnam}

\begin{itemize}[leftmargin=1.5em,itemsep=0em,topsep=0.1em]
    \item Researched few-shot continual learning for language models, designing
    data-augmentation and generalization-preserving training methods;
    equal-contribution first author at EMNLP 2024 and NAACL 2025
\end{itemize}

\end{rSection}

%----------------------------------------------------------------------------------------
% SELECTED PROJECTS AND OPEN SOURCE
%----------------------------------------------------------------------------------------

\begin{rSection}{Selected Projects and Open Source}

\textbf{\href{https://github.com/FSoft-AI4Code/CodeWiki}{CodeWiki}}
\hfill 1,500+ GitHub stars\\
Agentic framework for holistic documentation of large codebases; supports seven
languages and ships with an evaluation benchmark for repository-level understanding

\vspace{0.15em}

\textbf{\href{https://github.com/anhnh2002/XTTSv2-Finetuning-for-New-Languages}
{XTTSv2 Fine-Tuning for New Languages}}
\hfill 200+ GitHub stars\\
End-to-end toolkit for adapting the XTTSv2 speech-synthesis model to new languages,
covering data prep, tokenizer extension, training, and inference

\vspace{0.15em}

\textbf{AI for Real Action 2025, Champion (Team Converged)}
\hfill January 2026\\
First prize among nearly 400 teams in Vietnam's first national AI competition,
delivering a production-style AI solution under competition time constraints

\end{rSection}

%----------------------------------------------------------------------------------------
% PUBLICATIONS
%----------------------------------------------------------------------------------------

\begin{rSection}{Publications}

\begin{enumerate}[leftmargin=1.8em,label={[\arabic*]},itemsep=0.15em,topsep=0.1em]

    \item \textbf{CodeWiki: Evaluating AI's Ability to Generate Holistic
    Documentation for Large-Scale Codebases.}
    A. Nguyen Hoang, M. Le-Anh, B. Le, N. D. Q. Bui.
    \textit{Findings of ACL 2026.}
    \href{https://arxiv.org/abs/2510.24428}{[Paper]}

    \item \textbf{Beyond Repository Boundaries: Cross-Repository Graph
    Retrieval for Code Generation.}
    M. Le-Anh, N. Le Hai, Q. Tran, A. Nguyen Hoang, L. Ngo Van, B. Le,
    N. D. Q. Bui.
    \textit{Findings of EMNLP 2026.}
    % TODO: add \href{<arxiv-url>}{[Paper]} once the preprint is available

    \item \textbf{RustPrint: Documentation-Guided Agentic Codebase Migration
    from C to Rust.}
    M. Le-Anh, A. Nguyen Hoang, B. Le, N. D. Q. Bui.
    \textit{Preprint, 2026.}
    \href{https://arxiv.org/abs/2605.14634}{[Paper]}

    \item \textbf{Mutual-Pairing Data Augmentation for Few-Shot Continual
    Relation Extraction.}
    N. H. Anh*, Q. Tran*, T. X. Nguyen*, et al.
    \textit{NAACL 2025.}
    \href{https://aclanthology.org/2025.naacl-long.205/}{[Paper]}

    \item \textbf{Preserving Generalization of Language Models in Few-Shot
    Continual Relation Extraction.}
    Q. Tran*, N. X. Thanh*, N. H. Anh*, et al.
    \textit{EMNLP 2024.}
    \href{https://aclanthology.org/2024.emnlp-main.763/}{[Paper]}

\end{enumerate}
\textit{* equal contribution.}

\end{rSection}

%----------------------------------------------------------------------------------------
% EDUCATION
%----------------------------------------------------------------------------------------

\begin{rSection}{Education}

\textbf{Hanoi University of Science and Technology} \hfill 2020 -- 2024\\
Bachelor of Science in Computer Science, GPA 3.75/4.0
\hfill Ranked 15th of 296 (Top 5\%)

\vspace{0.15em}

\textbf{Techcombank Overseas PhD Study Sponsorship} (up to USD 333,000)
\hfill April 2026

\end{rSection}

%----------------------------------------------------------------------------------------
% LANGUAGES
%----------------------------------------------------------------------------------------

\begin{rSection}{Languages}

Vietnamese (native), English (professional working proficiency, IELTS 7.0)

\end{rSection}

\end{document}
